\section{Capture efficiencies from the diffusion tensor}
\label{app:woo}

The orientation factors $A_{h}(p)$ of \cref{eq:Afactors} follow from Woo's
treatment of a dislocation loop as an absorber in an anisotropically diffusing
medium~\citep{woo1981sink,woo1982intrinsic,woo1988theory}, and in the tensorial
form given by Po~\citep{po2026tensorialdad} the correction to the sink strength of a loop
of axis $\bm\alpha$ is
\begin{equation}
  \frac{\tr\Dten+\bm\alpha^{T}\Dten\bm\alpha}{4\overline D},
  \label{eq:tensorialDAD}
\end{equation}
which is universal in loop orientation and in the degree of anisotropy. For a
diagonal tensor with basal components $D_\perp$ and axial component
$D_\parallel$, and with $p=(D_\parallel/D_\perp)^{1/6}$ and
$\overline D=(\det\Dten)^{1/3}=D_\perp^{2/3}D_\parallel^{1/3}$,
\cref{eq:tensorialDAD} evaluates on a basal loop ($\bm\alpha\parallel[0001]$) to
$A_c(p)=p$ and on a prismatic loop ($\bm\alpha\perp[0001]$) to
$A_a(p)=(p+p^{-2})/2$.

The two forms differ in a way that has quantitative consequences and is easy to
overlook. $A_c$ is monotone increasing, so raising the anisotropy of a species
monotonically raises its capture by basal loops. $A_a$ has a \emph{minimum} at
$p=2^{1/3}=1.2599$, so raising the anisotropy from $p=1$ first \emph{lowers}
prismatic capture. A single change in $p^{v}$ therefore moves the two habits in
opposite directions. Setting $p^{v}=1$ where a fit had placed it at $1.179$
lowers $Z_{\mathrm{basal}}(v)$ by a factor of $0.848$ while \emph{raising}
$Z_{\mathrm{prism}}(v)$ by $1.053$ --- starving basal growth and accelerating
prismatic shrinkage at once. Any statement about the effect of anisotropy that
does not distinguish the habits is incomplete.

The co-growth condition of \cref{sec:coexistence} follows directly. Both growth
inequalities in \cref{eq:growconds} reduce to bounds on the single arrival ratio
$x$, and the window between them has width $g(p^{v})/g(p^{i})$ with
$g(p)=2/(1+p^{-3})$, strictly increasing --- so the window is open exactly when
$p^{i}<p^{v}$.

Two limits on this result should be stated. It governs the \emph{absorption}
flux only: cascade nucleation $G_k$ is added independently and no anisotropy
affects it, so the criterion does not govern net population evolution. And an
isotropic calculation is not a zero of the criterion but another parameter
point, so a ratio taken against one cannot test a statement about signs.

\section{The moment closure, and the three rules that keep it realizable}
\label{app:moments}

$\Delta_k=q^k_{sL}n^k_{sL}/(c^k_{sL})^{2}\ge1$ is the Cauchy--Schwarz inequality
on a non-negative measure. A state with $\Delta_k<1$ --- or with $q<0$ outright
--- is therefore not an inaccurate description of a distribution; it is not the
moment triple of any distribution. Keeping the triple realizable constrains how
each channel debits $q$, and three rules cover every channel in the model.

\begin{enumerate}[itemsep=3pt]
\item \textbf{Nucleation deposits at its declared size:} $J^k_{sL}
      (m^{\mathrm{nuc}}_k)^{2}$. This is genuine broadening,
      $J(m^{\mathrm{nuc}}-\bar m)^{2}\ge0$, and it is correct because the channel
      really does deposit at $m^{\mathrm{nuc}}_k$.
\item \textbf{The floor current debits at $m_{\min}$:}
      $-m_{\min}^{2}\Phi^k_{\min}$. This is a genuine statement about the lower
      tail.
\item \textbf{Channels with no distribution of their own are
      shape-preserving:} $\mathrm dq=\Delta_k(2\bar m\,\mathrm dc
      -\bar m^{2}\mathrm dn)$. Annealing and coalescence carry no size spectrum,
      so they must not be allowed to imply one.
\end{enumerate}

Rule 3 is where the sign is easily inverted, and the term is the largest
in the equation. A coalescing loop does not leave the family --- it
\emph{merges}. Its defects stay; only the loop--network content is lost; and
fewer loops holding the same content are \emph{larger} loops. Coalescence
therefore \emph{raises} $q$. Debiting it by the removed number times
$\langle m^{2}\rangle$, which is the natural-looking reading, inverts the sign of
the term that carries $98\%$ of the loop-number loss. The symptom in a
calculation is not an error but a slow drift: on one nine-family march, $q$
crossed zero at 1\,dpa and was negative on $40\%$ of the basal nodes at 10\,dpa.

The legacy annealing surrogate cannot be made consistent, only avoided.
It debits content at a \emph{fixed} size, so the pair $(\mathrm dn,\mathrm dc)$
it declares corresponds to a removed sub-population of mean $M\ne\bar m$.
Debiting $q$ proportionally over-removes; debiting at $M$ --- the smallest value
Cauchy--Schwarz permits for the removed part --- lowers $\Delta_k$ by
$a(M/\bar m-1)^{2}$ on a narrow family. No debit is realizable, because the
surrogate is removing loops of a size that is not there. That is an independent
argument for the peripheral-emission channel of \cref{app:emission}, which
deletes the fitted lifetimes outright.

The gated channels --- the small-size leak and the grain-boundary removal of
\cref{eq:gbloop} --- need no rule, because removing the mass of the distribution
on one side of a threshold leaves a genuine sub-measure and $\Delta_k\ge1$
survives by construction.

\section{Peripheral emission and the loop binding energy}
\label{app:emission}

The mean-field model represents loop annealing by two fitted lifetimes, one per
character. They are replaced here by an emission channel with no fitted rate, in
which a loop exchanges vacancies across its own perimeter against its own
equilibrium concentration.

The binding energy of a vacancy to a loop of $m$ defects is the difference of
the loop self-energies,
\begin{equation}
  E^{b}(m)=E^{f}_v+E_L(m-1)-E_L(m),
  \label{eq:Eb}
\end{equation}
with $E_L$ the anisotropic elastic self-energy of a circular loop of radius
$R=\lambda_k\sqrt m$ plus, for a faulted family, its stacking-fault term:
\begin{equation}
  E_L(m)=\frac{\mu\,(\bvec\!\cdot\!\nvec)_k^{2}\,R}{2(1-\nu)}
         \left[\ln\frac{8R}{r_c}-2\right]
       + \pi R^{2}\gamma_{\mathrm{sf}}\,\delta_{k,c_f}.
  \label{eq:EL}
\end{equation}
The local equilibrium vacancy concentration at the loop periphery is then
$c^{v,\mathrm{eq}}_k(m)=c^{v,\mathrm{eq}}\exp\bigl[-E^{b}(m)/k_BT\bigr]$, and the
emission current is the same capture expression run backwards against it,
\begin{equation}
  S^{k}_{vL}=\Pi^k_{sL}\,\overline D^{v}Z_{vk,v}\;
             c^{v,\mathrm{eq}}_k(\bar m^k_{sL}),
  \label{eq:Semit}
\end{equation}
which enters \cref{eq:master-c} through $\spvec S_{vL}$ and
\cref{eq:block-c} with the opposite sign.

Three properties recommend this over a fitted lifetime. \emph{One expression
serves both characters}: the same mechanism shrinks a vacancy loop, by emitting
the vacancies it holds, and grows an interstitial one, by emitting vacancies that
recombine with its interstitials, so no separate interstitial-loop annealing
model is needed. \emph{The climb threshold emerges rather than being imposed}:
growth requires the arriving flux to exceed \cref{eq:Semit}, which reproduces the
classical critical radius without it appearing as a parameter. And \emph{it
removes content, not loops}, which is physically correct --- emission shrinks a
loop continuously --- but has a consequence that must be recognized: after this
replacement, coalescence becomes the \emph{only} sink on vacancy loop number.
That is stated in \cref{sec:conclusions} as the leading open problem, and it is a
consequence of doing the emission physics correctly rather than an artifact of
it.

\section{State-vector layout and the conservation accumulators}
\label{app:state}

The integrated state at each node is a single flat vector whose width depends on
the model controls. The layout is fixed and every index keeps its meaning as the
controls add fields:

\begin{center}\small
\begin{tabular}{@{}llr@{}}
\toprule
Indices & Contents & Count \\
\midrule
$0$--$3$   & mobile: $c^{v},c^{i},c^{2i},c^{3i}$ & 4 \\
$4$--$7$   & number of families $0$--$3$ & 4 \\
$8$--$11$  & content of families $0$--$3$ & 4 \\
$12$--$17$ & accumulators: production, recombination, sink absorption,
             each per character & 6 \\
$18$       & network dislocation density $\rho_N$ & 1 \\
$19$--$23$ & number of families $4$--$8$ & 5 \\
$24$--$28$ & content of families $4$--$8$ & 5 \\
$29$--$37$ & second content $q$ of all nine families & 9 \\
$38$--$39$ & grain-boundary loop ledger, per character & 2 \\
\bottomrule
\end{tabular}
\end{center}

Widths are therefore 19 (four families, two moments), 29 (nine families, two
moments), 38 (nine families, three moments) and 40 with the boundary ledger.
Two properties of this layout are worth recording because both were established
by their failure modes.

The appended families are appended, not interleaved. Every index below
19 means what it meant in the four-family model, so a state written by one
configuration is readable by another up to its own width. The corollary is that
a narrower state is always \emph{accepted} by a wider solver, the appended slots
defaulting to zero --- which is exactly why a mismatch between the width a driver
allocates and the width the model requires is silent rather than fatal, and why
the checkpoint fingerprint records the actual width.

The block written to the spatial code is ordered the other way. It is
moment-major --- all nine numbers, then all nine contents, then all nine $q$ ---
so the bridge between the two is a scatter and not a copy. Exactly one routine
knows that mapping; performing it at a call site is what once put an 18-wide
block into an 8-wide slice.

The accumulators are per-interval; the ledger is cumulative. The six
conservation accumulators are reset at the start of each dose interval, so a
stored snapshot holds one interval's production and must be summed to give a
running total. The two grain-boundary accumulators are never reset and are
already cumulative. Differencing the wrong one against the wrong other gives a
channel that appears to carry $6600$ times total production, or none at all,
depending on which way the error runs.
